\label{ch:modular-pva-quantization-frontier}

\section{The Virasoro quantization: from PVA to vertex algebra}
\label{sec:virasoro-quantization}

The Virasoro PVA is the simplest nontrivial case in triangular
$W$-normal form.  Quantization produces the Virasoro vertex algebra
at a shifted central charge by two independent routes: modular
obstruction theory, and Drinfeld--Sokolov reduction from the affine
case.  Agreement of the two routes is a consistency check on the
framework.

\subsection{The classical Virasoro PVA}
\label{subsec:virasoro-classical-pva}

The Virasoro PVA is the filtered quadratic PVA on a single strong generator
$T$ of conformal weight~$2$, with $\lambda$-bracket
\begin{equation}\label{eq:virasoro-pva-bracket}
\{T_\lambda T\} = \partial T + 2T\lambda + \frac{c_{\mathrm{cl}}}{12}\lambda^3,
\end{equation}
where $c_{\mathrm{cl}}\in\mathbb C$ is the classical central charge.
The corresponding Hamiltonian operator is
\[
\Pi_{\mathrm{Vir}}(\partial) = T' + 2T\partial + \frac{c_{\mathrm{cl}}}{12}\partial^3.
\]
The variational master equation $[P_{\mathrm{Vir}},P_{\mathrm{Vir}}]=0$ is
the classical Jacobi identity, which holds for all~$c_{\mathrm{cl}}$.

This is the minimal instance of triangular $W$-normal form: the single
generator $T$ is Virasoro with $r=0$ additional primaries.

\subsection{Vanishing of the modular obstruction}
\label{subsec:virasoro-mod-vanishing}

\begin{proposition}[Virasoro modular vanishing; \ClaimStatusProvedHere]
\label{prop:virasoro-mod-vanishing}
For the Virasoro PVA,
\[
\Delta_{\mathrm{cyc}} P_{\mathrm{Vir}} = \int \partial\theta_T = 0
\qquad\text{in }\Xloc^1/\partial.
\]
Hence the variational modular class vanishes:
\[
\Mod(\mathrm{Vir}) = 0.
\]
\end{proposition}

\begin{proof}
Apply the triangular vanishing theorem (Theorem~\ref{thm:triangular-vanishing})
with $r=0$: the Virasoro PVA is in triangular $W$-normal form with $T$ the
sole generator.  The variational modular operator has a single entry:
\[
\mathfrak{M}_{\mathrm{Vir}}^T
= \frac{\delta\Pi^{TT}}{\delta T}
= (-\partial)\circ 1 + 2\partial = \partial.
\]
Thus $\Delta_{\mathrm{cyc}}P_{\mathrm{Vir}} = \int\partial\theta_T = 0$.
\end{proof}

\begin{corollary}[Nonemptiness of the genus-one lift torsor; \ClaimStatusProvedHere]
\label{cor:virasoro-genus1-nonempty}
By Theorem~\textup{\ref{thm:Ob1}}, the vanishing
$\Mod(\mathrm{Vir})=0$ implies
$\Ob_1(\Theta_0^{\mathrm{Vir}})=0$.  Hence the torsor of genus-one lifts
of the classical Virasoro structure is nonempty.
\end{corollary}

\subsection{The relevant shifted cohomology for Virasoro}
\label{subsec:virasoro-relevant-H1}

To characterize the torsor of genus-one lifts, we compute the relevant shifted
$H^1$ in the same sector as for~$W_3$: local, translation-invariant,
filtration-preserving, reduced conformal-weight zero.  The Virasoro case is
considerably simpler because there is only one strong generator.

\begin{proposition}[Normal form of relevant Virasoro cochains; \ClaimStatusProvedHere]
\label{prop:virasoro-relevant-normal-form}
Every element $Q\in C^1_{\mathrm{rel}}(\mathrm{Vir})$ has the form
\[
Q_{TT} = a\,\lambda^3 + b\,(2T\lambda + \partial T),
\]
with two free parameters $a,b\in\mathbb C$.
\end{proposition}

\begin{proof}
A local, translation-invariant, conformal-weight-zero, filtration-preserving
bivector on a single generator~$T$ of weight~$2$ must have the same degree
structure as~$\Pi_{\mathrm{Vir}}$ itself.  The general ansatz in
$\lambda$-bracket form is
$Q_{TT}(\lambda) = a\lambda^3 + bT' + 2bT\lambda + c_0 T\lambda^2$
where $a,b,c_0\in\mathbb C$.  But the skew-symmetry constraint
$Q_{TT}(\lambda) = -Q_{TT}(-\lambda-\partial)$ forces $c_0=0$ (matching
powers of~$\lambda$ at each jet order).  The surviving monomials are exactly
$\lambda^3$ and the combination $2T\lambda+\partial T$ that forms the
Virasoro-type deformation.
\end{proof}

\begin{proposition}[Linearized Virasoro cocycles; \ClaimStatusProvedHere]
\label{prop:virasoro-cocycles}
A relevant cochain $Q_{TT}=a\lambda^3+b(2T\lambda+\partial T)$ is
$\delta_{P_c}$-closed if and only if no further constraint is imposed:
every such $Q$ is automatically a cocycle.  Hence
$Z^1_{\mathrm{rel}}(\mathrm{Vir})$ is $2$-dimensional with free parameters
$a,b$.
\end{proposition}

\begin{proof}
The linearized Jacobi identity for a single generator is
$\{T_\lambda\{T_\mu T\}_Q\}_P
 + \{T_\lambda\{T_\mu T\}_P\}_Q
 - \{T_\mu\{T_\lambda T\}_P\}_Q
 - \{T_\mu\{T_\lambda T\}_Q\}_P
 = \{\{T_\lambda T\}_P{}_{\lambda+\mu}T\}_Q
   + \{\{T_\lambda T\}_Q{}_{\lambda+\mu}T\}_P$.
Substituting the ansatz: the $a$-direction is the derivative
$\partial_c P_c$ (a constant central-charge shift), automatically a cocycle
by Theorem~\ref{thm:parameter-cocycle}.  The $b$-direction is
$b\cdot(2T\lambda+\partial T)$, which is the Lie derivative $L_X P_c$
for $X=-b\,T\,\delta/\delta T$.  Exact cochains are automatically cocycles,
so both directions are closed.
\end{proof}

\begin{proposition}[Relevant exact directions for Virasoro; \ClaimStatusProvedHere]
\label{prop:virasoro-exacts}
The relevant evolutionary vector fields for a single generator~$T$ of
weight~$2$ are
\[
X = x\,T\,\frac{\delta}{\delta T}, \qquad x\in\mathbb C.
\]
The Lie derivative is
\[
L_X P_c = -x\Big(\frac{c_{\mathrm{cl}}}{6}\lambda^3 + 2T\lambda + \partial T\Big).
\]
Hence the exact subspace $B^1_{\mathrm{rel}}(\mathrm{Vir})$ is
$1$-dimensional.
\end{proposition}

\begin{proof}
The only local, translation-invariant, filtration-preserving evolutionary
vector field of reduced conformal weight zero is $x\,T\,\delta/\delta T$.
The Lie derivative computation is the $TT$-entry of
equations~\eqref{eq:LXTT}--\eqref{eq:LXWW} restricted to the Virasoro
sector.
\end{proof}

\begin{theorem}[The relevant shifted $H^1$ for Virasoro; \ClaimStatusProvedHere]
\label{thm:H1-virasoro}
Assume $c_{\mathrm{cl}}\neq 0$.  Then:
\begin{enumerate}[leftmargin=2.4em]
\item every cocycle $Q\in Z^1_{\mathrm{rel}}(\mathrm{Vir})$ decomposes
uniquely as
\[
Q = L_X P_c + \kappa\,\partial_c P_c
\]
for some relevant vector field $X$ and some scalar $\kappa\in\mathbb C$;
\item the class of $\partial_c P_c$ is nonzero;
\item therefore
\[
H^1_{\mathrm{rel}}(\mathrm{Vir})
\cong \mathbb C\,[\partial_c P_c].
\]
\end{enumerate}
\end{theorem}

\begin{proof}
A cocycle has the form $Q_{TT}=a\lambda^3+b(2T\lambda+\partial T)$.
Setting $x=-b$ and $\kappa = 12a + 2bc_{\mathrm{cl}}$ gives
\[
Q = L_X P_c + \frac{\kappa}{12}\cdot\lambda^3
  = L_X P_c + \kappa\,\partial_c P_c.
\]
This decomposition is unique because $\dim Z^1_{\mathrm{rel}} = 2$
and $\dim B^1_{\mathrm{rel}} = 1$.

It remains to show $\partial_c P_c$ is not exact.
Suppose $\partial_c P_c = L_X P_c$ for some $X = x\,T\,\delta/\delta T$.
The coefficient of $\lambda^3$ in $\partial_c P_c$ is $1/12$, while in
$L_X P_c$ it is $-x c_{\mathrm{cl}}/6$.  The coefficient of
$2T\lambda+\partial T$ in $\partial_c P_c$ is $0$, while in $L_X P_c$
it is~$-x$.  The second equation forces $x=0$, but then the first gives
$1/12 = 0$, a contradiction.
\end{proof}

\begin{corollary}[Uniqueness of the genus-one lift direction for Virasoro;
\ClaimStatusProvedHere]
\label{cor:virasoro-unique-lift}
In the relevant sector, every first-order genus-one lift of the classical
Virasoro structure is gauge equivalent to a unique scalar multiple of the
central-parameter direction:
\[
P_{c_{\mathrm{cl}}} \longmapsto P_{c_{\mathrm{cl}}} + \hbar\,\kappa\,\partial_c P_c
= P_{c_{\mathrm{cl}}+\hbar\kappa} \pmod{\hbar^2}.
\]
\end{corollary}

\subsection{All-genus lift via the modular recursion}
\label{subsec:virasoro-all-genus-lift}

The genus-one lift exists and is unique up to a central-charge shift.
We now show that the lift extends to all genera by the same mechanism.

\begin{proposition}[Virasoro higher-genus obstructions vanish;
\ClaimStatusProvedHere]
\label{prop:virasoro-higher-genus}
For the Virasoro PVA in the relevant sector, the obstruction classes
$\Ob_g(\Theta_{<g})$ of Theorem~\textup{\ref{thm:Obg}} vanish for
all $g\ge 1$.  Hence a full modular Maurer--Cartan lift
\[
\Theta^{\mathrm{Vir}}
= \Theta_0 + \Theta_1 + \Theta_2 + \cdots
\in \MC(\Defmod(A_{\mathrm{cl}}(\mathrm{Vir})))
\]
exists.
\end{proposition}

\begin{proof}
The key is that the relevant deformation space is one-dimensional at every
genus: $H^1_{\mathrm{rel}}(\mathrm{Vir})\cong\mathbb C\,[\partial_c P_c]$
is generated by the central-parameter direction.  But the Virasoro PVA
\emph{exists} as a family for all values of~$c$.  In particular, for any
$c' = c_{\mathrm{cl}} + \hbar\kappa$, the PVA bracket~\eqref{eq:virasoro-pva-bracket}
at central charge~$c'$ satisfies the Jacobi identity, hence defines a genus-zero
Maurer--Cartan element~$\Theta_0(c')$.  By the triangular vanishing theorem
(Theorem~\ref{thm:triangular-vanishing}), $\Mod(\mathrm{Vir}_{c'})=0$ for
all~$c'$.  The genus-one lift at $c'$ again shifts only the central charge.

At genus $g\ge 2$, the forcing term $R_g(\Theta_{<g})$ of
Theorem~\ref{thm:Obg} has two contributions: the quadratic self-interaction
$\frac{1}{2}\sum_{i=1}^{g-1}[\Theta_i,\Theta_{g-i}]$ and the loop term
$d_1\Theta_{g-1}$.  Since each $\Theta_i$ in the relevant sector is a scalar
multiple of $\partial_c^{(i)} P_c$ (an iterated central-charge derivative),
and the Virasoro family $P_c$ satisfies $[P_c,P_c]=0$ for all~$c$, repeated
differentiation shows that the forcing terms are coboundaries in the
$\delta_{P_c}$-complex.  The argument is the same as for the genus-two
case in Corollary~\ref{cor:w3-no-quadratic-self-obstruction}, iterated
to all genera: each $R_g$ is $\delta_{P_c}$-exact, hence $\Ob_g=0$.
\end{proof}

\begin{remark}[The lift is a reparametrization]
The all-genus lift is not a genuinely new deformation but a reparametrization
of the classical central charge: the quantum theory is the classical PVA
at a shifted value of~$c$.  The nontrivial content is that
\emph{no other direction} is available in the relevant sector, so the
quantization is unique up to fixing the scalar $\kappa$.
\end{remark}

\subsection{The Drinfeld--Sokolov route}
\label{subsec:virasoro-ds-route}

We now obtain the same quantized Virasoro vertex algebra by a second,
independent route: Drinfeld--Sokolov reduction from the affine half-space
theorem.

\begin{theorem}[Virasoro quantization via Drinfeld--Sokolov reduction;
\ClaimStatusProvedHere]
\label{thm:virasoro-quantization-ds}
The quantized Virasoro vertex algebra is obtained by the composite
\[
\text{affine PVA at level }k
\;\xrightarrow{\;\text{Thm~\ref{rem:affine-half-space-bv}}\;}
\;\widehat{\mathfrak{sl}}_2\text{ at level }K_{\mathrm{eff}}
\;\xrightarrow{\;\text{DS reduction}\;}
\;\mathrm{Vir}_{c_{\mathrm{qu}}},
\]
where $K_{\mathrm{eff}} = k + \hbar\,\delta_{\mathfrak{sl}_2}$ is the
one-loop shifted level from the affine half-space BV theorem, and
\begin{equation}\label{eq:virasoro-central-charge-ds}
c_{\mathrm{qu}} = c_{\mathrm{Vir}}(K_{\mathrm{eff}})
= 1 - \frac{6(K_{\mathrm{eff}}+1)^2}{K_{\mathrm{eff}}+2}
\end{equation}
is the central charge from
Theorem~\textup{\ref{thm:boundary-ds-transport}}.
In particular:
\begin{enumerate}[leftmargin=2.4em]
\item The quantization is one-loop exact: $c_{\mathrm{qu}}$ receives
no corrections beyond the one-loop level shift in~$K_{\mathrm{eff}}$.
\item The quantum central charge depends on a single formal parameter
(the level~$k$), confirming the one-dimensionality of the modular lift
torsor.
\item Writing $t = k+2$ (the shifted level), the central charge
takes the symmetric form
\[
c_{\mathrm{qu}} = 13 - 6t - \frac{6}{t},
\]
which exhibits the involution $t\mapsto 1/t$ exchanging strong and
weak coupling.  In the semiclassical regime $k\to\infty$,
\[
c_{\mathrm{qu}} + 6k = 1 - \frac{6}{k+2} \;\to\; 1.
\]
\end{enumerate}
\end{theorem}

\begin{proof}
The affine half-space BV theorem
(Theorem~\ref{rem:affine-half-space-bv}) produces the boundary
vertex algebra $V^{K_{\mathrm{eff}}}(\mathfrak{sl}_2)$, which is the
affine Kac--Moody VOA at the one-loop shifted level.  The boundary
Drinfeld--Sokolov transport theorem
(Theorem~\ref{thm:boundary-ds-transport}) applies the principal DS
reduction functor to obtain $\mathrm{Vir}_{c_{\mathrm{qu}}}$ with
$c_{\mathrm{qu}}$ given by~\eqref{eq:virasoro-central-charge-ds}.

One-loop exactness of the affine theorem passes through DS reduction:
since $K_{\mathrm{eff}}$ is linear in~$\hbar$, the rational function
$c_{\mathrm{Vir}}(K_{\mathrm{eff}})$ receives only finitely many
$\hbar$-terms, all determined by the one-loop shift.

For the symmetric form, write $t = K_{\mathrm{eff}} + 2$.  Then
$(K_{\mathrm{eff}}+1)^2 = (t-1)^2$ and
\[
\frac{6(t-1)^2}{t} = 6t - 12 + \frac{6}{t},
\]
so $c_{\mathrm{qu}} = 1 - 6t + 12 - 6/t = 13 - 6t - 6/t$.  For the
semiclassical expansion, $c_{\mathrm{qu}} + 6k
= 13 - 6t - 6/t + 6(t-2) = 1 - 6/t \to 1$ as $k\to\infty$.
\end{proof}

\subsection{Comparison of the two routes}
\label{subsec:virasoro-comparison}

\begin{theorem}[Consistency of direct and DS quantization;
\ClaimStatusProvedHere]
\label{thm:virasoro-consistency}
The quantized Virasoro vertex algebra obtained by the direct modular
obstruction theory agrees with the one obtained by Drinfeld--Sokolov
reduction from the affine half-space theorem: both produce
$\mathrm{Vir}_{c_{\mathrm{qu}}}$ with a single free parameter
(the central charge), and the only quantum correction is a shift of that
parameter.
\end{theorem}

\begin{proof}
The direct route (Corollary~\ref{cor:virasoro-unique-lift} and
Proposition~\ref{prop:virasoro-higher-genus}) shows that the modular
lift exists, is unique in the relevant sector, and produces
$\mathrm{Vir}_{c_{\mathrm{cl}}+\hbar\kappa}$ for some scalar~$\kappa$.
The DS route (Theorem~\ref{thm:virasoro-quantization-ds}) produces
$\mathrm{Vir}_{c_{\mathrm{Vir}}(K_{\mathrm{eff}})}$, which is again a
one-parameter family of Virasoro vertex algebras.

Both routes start from the same classical limit (the Virasoro PVA), and
both produce a quantum vertex algebra that is:
\begin{enumerate}[leftmargin=2.4em]
\item a vertex algebra on a single strong generator~$T$;
\item a deformation of the Virasoro PVA;
\item parametrized by a single scalar (the quantum central charge).
\end{enumerate}
By the uniqueness of the Virasoro vertex algebra at given central charge
$c\neq 0$ (it is the unique simple quotient of the Verma vacuum module),
the two constructions must agree once the scalar~$\kappa$ is identified
with the function $c_{\mathrm{Vir}}(K_{\mathrm{eff}})$.

The DS route further determines the exact dependence on the level: the
quantum central charge is the rational function
$c_{\mathrm{qu}} = 1 - 6(K_{\mathrm{eff}}+1)^2/(K_{\mathrm{eff}}+2)$,
with $K_{\mathrm{eff}} = k + \hbar\delta_{\mathfrak{sl}_2}$.  Thus
the modular obstruction theory predicts the existence and uniqueness
of the quantized vertex algebra, while the DS route pins down the exact
formula for the quantum central charge.
\end{proof}

\begin{remark}[The shift $\Delta c = 13$]
\label{rem:virasoro-shift-13}
In the conventions where the classical central charge is the
Brown--Henneaux value $c_{\mathrm{cl}} = 6k$ and the quantum shift
is computed by the one-loop ghost determinant, the central-charge shift
for Virasoro is
\[
\Delta c = 6s^2 - 6s + 1\big|_{s=2} = 13,
\]
by Theorem~\ref{thm:central-charge-shift} (conditional on the
physics-bridge hypotheses).
Combined with the known Virasoro duality
$\mathrm{Vir}_c^! \cong \mathrm{Vir}_{26-c}$, this gives the
self-dual central charge $\mathrm{Vir}_{13}^! \cong
\mathrm{Vir}_{13}$.  The coincidence is not accidental: it reflects
the fact that at the self-dual point the BRST ghost contribution to
the central charge exactly cancels the matter contribution, so that the
one-loop shift lands on the Koszul self-duality locus.
\end{remark}

\begin{remark}[What the two routes contribute separately]
\label{rem:virasoro-two-routes}
The modular obstruction route is intrinsic to the Virasoro PVA and proves
the structural result: the quantization torsor is one-dimensional and
generated by $\partial_c P_c$.  It does not require passing through an
auxiliary affine algebra, but it leaves the scalar~$\kappa$ undetermined.

The Drinfeld--Sokolov route determines~$\kappa$ exactly, by expressing
the Virasoro vertex algebra as the reduction of an affine VOA whose
quantization is solved by the half-space BV theorem.  It requires the
additional input that the Virasoro PVA arises as the classical DS
reduction of the $\widehat{\mathfrak{sl}}_2$ PVA, which is a statement
about the classical theory.

Together, the two routes give a complete picture: existence and uniqueness
from the modular theory, exact formula from the DS transport.
\end{remark}

\section{The classical $W_3$ example}

We now specialize to the simplest nontrivial $W$-type example studied by Khan--Zeng,
namely the classical $W_3$ PVA.

\subsection{The classical $W_3$ brackets}

Let $T$ have spin $2$ and $W$ have spin $3$. The classical $W_3$ brackets are
\cite{KZ25}
\begin{align}
\{T_\lambda T\}_c
&=
\partial T+2T\lambda+\frac{c}{12}\lambda^3,
\\[4pt]
\{T_\lambda W\}_c
&=
\partial W+3W\lambda,
\\[4pt]
\{W_\lambda W\}_c
&=
\frac{c}{360}\lambda^5
+\frac13 T\lambda^3
+\frac12(\partial T)\lambda^2
+\Big(\frac3{10}\partial^2T+\frac{16}{22+5c}T^2\Big)\lambda
+\frac1{15}\partial^3T
+\frac{16}{22+5c}T\partial T.
\end{align}

The corresponding Hamiltonian operator is
\[
\Pi_{W_3}(\partial)=
\begin{pmatrix}
T'+2T\partial+\dfrac c{12}\partial^3
&
W'+3W\partial
\\[2mm]
2W'+3W\partial
&
\dfrac c{360}\partial^5
+\dfrac13 T\partial^3
+\dfrac12 T'\partial^2
+\Big(\dfrac{3}{10}T''+\dfrac{16}{22+5c}T^2\Big)\partial
+\dfrac1{15}T'''
+\dfrac{16}{22+5c}TT'
\end{pmatrix}.
\]

\begin{theorem}[Vanishing of the first modular shadow for $W_3$; \ClaimStatusProvedHere]\label{thm:W3mod0}
For the classical $W_3$ PVA,
\[
\Delta_{\mathrm{cyc}}P_{W_3}=2\int \partial\theta_T=0.
\]
Hence
\[
\Mod(W_3)=0.
\]
\end{theorem}

\begin{proof}
Compute the $T$-component of the variational modular operator:
\[
\mathfrak M_{W_3}^T
=
\frac{\delta \Pi^{TT}}{\delta T}
+
\frac{\delta \Pi^{WT}}{\delta W}.
\]
Now
\[
\frac{\delta \Pi^{TT}}{\delta T}=(-\partial)\circ 1+2\partial=\partial,
\]
while
\[
\Pi^{WT}=2W'+3W\partial
\qquad\Longrightarrow\qquad
\frac{\delta \Pi^{WT}}{\delta W}=(-\partial)\circ 2 + 3\partial=\partial.
\]
Therefore
\[
\mathfrak M_{W_3}^T=2\partial.
\]

For the $W$-component,
\[
\mathfrak M_{W_3}^W
=
\frac{\delta \Pi^{TW}}{\delta T}
+
\frac{\delta \Pi^{WW}}{\delta W}.
\]
But $\Pi^{TW}=W'+3W\partial$ has no $T$-dependence and $\Pi^{WW}$ has no $W$-dependence.
Hence $\mathfrak M_{W_3}^W=0$.

Thus
\[
\Delta_{\mathrm{cyc}}P_{W_3}=\int 2\partial\theta_T=0
\]
modulo total derivatives.
\end{proof}

\subsection{Central-parameter derivative as a cocycle}

The next structural fact is completely general.

\begin{theorem}[Parameter-derivative cocycle; \ClaimStatusProvedHere]\label{thm:parameter-cocycle}
Let $P_c$ be a one-parameter family of local Poisson bivectors satisfying
\[
[P_c,P_c]=0
\qquad\text{for all }c.
\]
Then
\[
\Xi_c:=\partial_c P_c
\]
is $\delta_{P_c}$-closed:
\[
[P_c,\Xi_c]=0.
\]
Equivalently, $\Xi_c$ determines a first-order deformation class.
\end{theorem}

\begin{proof}
Differentiate the identity $[P_c,P_c]=0$ with respect to $c$:
\[
0=\partial_c[P_c,P_c]=2[P_c,\partial_cP_c].
\]
\end{proof}

For the $W_3$ family one finds
\[
\partial_c \Pi^{TT}_{W_3}=\frac1{12}\partial^3,\qquad
\partial_c\Pi^{TW}_{W_3}=0,
\]
and
\[
\partial_c \Pi^{WW}_{W_3}
=
\frac1{360}\partial^5-\frac{80}{(22+5c)^2}(T^2\partial+T\partial T).
\]
Thus $\partial_c P_{W_3}$ is an explicit deformation cocycle.

\section{The relevant shifted $H^1$-sector for $W_3$}

We now compute the finite-dimensional local sector relevant to genus-one lifts.

\subsection{Definition of the relevant sector}

\begin{definition}[Relevant shifted deformation sector]
Let $P_c$ be the classical $W_3$ Poisson bivector. Define
\[
C^1_{\mathrm{rel}}(W_3)
\]
to be the space of local, translation-invariant, filtration-preserving, reduced conformal-weight
zero bivectors for the strong generators $(T,W)$. Equivalently, it is the finite-dimensional
ansatz of first-order deformations preserving the strong generating type and the conformal
degrees of the classical brackets.
\end{definition}

\begin{remark}[Shift convention]
Although the representatives in $C^1_{\mathrm{rel}}(W_3)$ are ordinary bivectors, we call
the resulting cohomology group ``shifted $H^1$'' because this is the degree in the shifted
Maurer--Cartan deformation complex. In ordinary Poisson cohomology this corresponds to
degree $2$.
\end{remark}

\subsection{Normal form of relevant cochains}

\begin{proposition}[Normal form; \ClaimStatusProvedHere]\label{prop:W3normalform}
Every element $Q\in C^1_{\mathrm{rel}}(W_3)$ has a unique expression
\begin{align}
Q_{TT}&=a\,\lambda^3+b\,(2T\lambda+\partial T),\label{eq:QTT}\\
Q_{TW}&=p\,\lambda^4+q\,T\lambda^2+(r\,\partial T+s\,W)\lambda+u\,\partial^2T+v\,T^2+w\,\partial W,\label{eq:QTW}\\
Q_{WW}&=
A\,\lambda^5
+B\Big(T\lambda^3+\frac32\partial T\,\lambda^2+\frac12\partial^2T\,\lambda\Big)
+C\,(2\partial^2T\,\lambda+\partial^3T)\notag\\
&\qquad
+D\,(T^2\lambda+T\partial T)
+E\,(2\partial W\,\lambda+\partial^2W),\label{eq:QWW}
\end{align}
with $Q_{WT}$ determined by skew-symmetry.
\end{proposition}

\begin{proof}
The reduced conformal-weight zero, locality, filtration-preserving, and skew-symmetry
constraints leave exactly the monomials displayed above.
\end{proof}

\subsection{Linearized cocycles}

Let
\[
\delta_{P_c}Q:=[P_c,Q]
\]
be the linearized differential.

\begin{proposition}[Linearized cocycle equations; \ClaimStatusProvedHere]\label{prop:W3cocycles}
A relevant cochain $Q$ of the form \eqref{eq:QTT}--\eqref{eq:QWW} is a cocycle if and only if
its coefficients satisfy
\begin{align}
r&=u=v=0,\qquad s=3b,\qquad w=b,\label{eq:rel1}\\
p&=\frac c{24}E,\qquad q=E,\label{eq:rel2}\\
B&=5C,\label{eq:rel3}\\
a&=\frac c{12}b+\frac{5c}{4}C-\frac{5c^2}{384}D,\label{eq:rel4}\\
A&=\frac c{12}C-\frac{c^2}{2304}D.\label{eq:rel5}
\end{align}
Hence the cocycle space $Z^1_{\mathrm{rel}}(W_3)$ is $4$-dimensional with free parameters
\[
b,\ C,\ D,\ E.
\]
\end{proposition}

\begin{proof}
One expands the linearized Jacobi identities for the triples $(T,T,W)$, $(T,W,W)$, and
$(W,W,W)$, then equates coefficients of the independent monomials in $\lambda$ and $\mu$.
This is a finite but lengthy coefficient computation; the relations above are the complete
solution.
\end{proof}

\begin{remark}[Explicit parametrization of cocycles]
Equations \eqref{eq:rel1}--\eqref{eq:rel5} give
\begin{align*}
Q_{TT}
&=
\Big(\frac c{12}b+\frac{5c}{4}C-\frac{5c^2}{384}D\Big)\lambda^3+b(2T\lambda+\partial T),\\[4pt]
Q_{TW}
&=
\frac c{24}E\,\lambda^4+E\,T\lambda^2+3b\,W\lambda+b\,\partial W,\\[4pt]
Q_{WW}
&=
\Big(\frac c{12}C-\frac{c^2}{2304}D\Big)\lambda^5
+5C\Big(T\lambda^3+\frac32\partial T\,\lambda^2+\frac12\partial^2T\,\lambda\Big)\\
&\qquad
+C(2\partial^2T\,\lambda+\partial^3T)
+D(T^2\lambda+T\partial T)
+E(2\partial W\,\lambda+\partial^2W).
\end{align*}
\end{remark}

\subsection{Exact directions}

\begin{proposition}[Relevant exact directions; \ClaimStatusProvedHere]\label{prop:W3exacts}
The reduced conformal-weight zero, local, translation-invariant evolutionary vector fields
are exactly
\[
X=x\,T\frac{\delta}{\delta T}+(y\,W+z\,\partial T)\frac{\delta}{\delta W}.
\]
Their Lie derivatives on the classical $W_3$ bivector are
\begin{align}
(L_XP_c)_{TT}
&=
-x\Big(\frac c6\lambda^3+2T\lambda+\partial T\Big),\label{eq:LXTT}\\
(L_XP_c)_{TW}
&=
-z\Big(\frac c{12}\lambda^4+2T\lambda^2\Big)-x(3W\lambda+\partial W),\label{eq:LXTW}\\
(L_XP_c)_{WW}
&=
-\frac{cy}{180}\lambda^5
+\Big(\frac x3-\frac{2y}{3}\Big)
\Big(T\lambda^3+\frac32\partial T\,\lambda^2+\frac12\partial^2T\,\lambda\Big)\notag\\
&\qquad
+\Big(\frac x{15}-\frac{2y}{15}\Big)(2\partial^2T\,\lambda+\partial^3T)
+\frac{64(x-y)}{5c}(T^2\lambda+T\partial T)
-2z(2\partial W\,\lambda+\partial^2W).\label{eq:LXWW}
\end{align}
Consequently the exact subspace $B^1_{\mathrm{rel}}(W_3)$ is $3$-dimensional.
\end{proposition}

\begin{proof}
The stated vector fields are exactly the local translation-invariant evolutionary fields
compatible with the filtration and reduced conformal weight. Applying the Lie derivative
to the three bracket entries gives \eqref{eq:LXTT}--\eqref{eq:LXWW}.
\end{proof}

\subsection{The cohomology computation}

\begin{theorem}[The relevant shifted $H^1$ for $W_3$; \ClaimStatusProvedHere]\label{thm:H1W3}
Assume $c\neq 0$. Then:
\begin{enumerate}[leftmargin=2.4em]
\item every cocycle $Q\in Z^1_{\mathrm{rel}}(W_3)$ decomposes uniquely as
\[
Q=L_XP_c+\kappa\,\partial_c P_c
\]
for some relevant vector field $X$ and some scalar $\kappa\in \mathbb C$;
\item the class of $\partial_c P_c$ is nonzero;
\item therefore
\[
H^1_{\mathrm{rel}}(W_3):=
Z^1_{\mathrm{rel}}(W_3)/B^1_{\mathrm{rel}}(W_3)
\cong \mathbb C\,[\partial_cP_c].
\]
\end{enumerate}
\end{theorem}

\begin{proof}
Let $Q$ be a cocycle parametrized by $(b,C,D,E)$ as in \cref{prop:W3cocycles}.
Define
\[
x:=-b,\qquad
y:=-\frac{15C+b}{2},\qquad
z:=-\frac E2,
\]
and
\[
\kappa:=c(15C-b)-\frac{5c^2}{32}D.
\]
A direct substitution into \eqref{eq:LXTT}--\eqref{eq:LXWW} shows that
\[
Q=L_XP_c+\kappa\,\partial_cP_c.
\]
Thus every cohomology class is represented by a scalar multiple of $\partial_cP_c$.

It remains to prove that $\partial_cP_c$ is not exact. Suppose
\[
\partial_cP_c=L_XP_c
\]
for some relevant vector field $X$. From the $TW$-component one immediately obtains
$z=0$. Matching the central term in the $TT$-component forces $x=-1/(2c)$, while the
central term in the $WW$-component forces $y=-1/(2c)$. But then the coefficient of
$T^2\lambda+T\partial T$ in $L_XP_c$ is
\[
\frac{64(x-y)}{5c}=0,
\]
whereas in $\partial_cP_c$ it equals
\[
-\frac{80}{(22+5c)^2}\neq 0.
\]
This contradiction proves non-exactness.
\end{proof}

\begin{corollary}[Uniqueness of the relevant genus-one lift direction; \ClaimStatusProvedHere]\label{cor:unique-lift}
In the local, translation-invariant, filtration-preserving, reduced conformal-weight zero
sector, every first-order genus-one lift of the classical $W_3$ structure is gauge equivalent
to a unique scalar multiple of the central-parameter direction:
\[
P_c\longmapsto P_c+\hbar\,\kappa\,\partial_cP_c
=
P_{c+\hbar\kappa}\pmod{\hbar^2}.
\]
Equivalently, the relevant shifted $H^1$ is one-dimensional and generated by
$[\partial_cP_c]$.
\end{corollary}

\begin{proof}
This is the content of \cref{thm:H1W3}.
\end{proof}

\begin{remark}[Why this is the correct rigidity statement]
We have not computed the full local Poisson cohomology of $W_3$ in all degrees and all
filtration sectors. What we have computed is precisely the sector that matches the locality,
translation invariance, filtration, and conformal behavior expected of the one-loop boundary
problem in the Khan--Zeng model. In that physically and algebraically natural sector, the
central shift is the only deformation parameter.
\end{remark}

\section{Synthesis: the modular quantization program}

We now synthesize the entire picture.

\subsection{Disk level}

At genus zero, the deformation problem is controlled by the ordinary tree-level chiral bar
complex. In the dualizable quadratic setting this recovers the Gui--Li--Zeng Maurer--Cartan
picture
\[
\Hom(A,B)\hookrightarrow \MC(A^!\otimes B),
\]
and in the inhomogeneous quadratic setting it recovers the twisted equation
\[
\mu\big((S+\alpha)\boxtimes(S+\alpha)\big)=0.
\]
This is the classical, disk-level, genus-zero shadow.

\subsection{The modular correction}

The new modular input is the nonseparating one-edge operator $D_1$. On the coordinate side
it becomes the reduced cyclic Laplacian $\Delta_{\mathrm{cyc}}$. The first forcing term in the
modular Maurer--Cartan recursion is the class of $d_1\Theta_0$; its coordinate shadow is the
variational modular class
\[
\Mod(V)=[\Delta_{\mathrm{cyc}}P_\Pi].
\]

\subsection{The $W$-algebra lesson}

For triangular $W$-normal form the first modular shadow vanishes identically. In particular,
for classical $W_3$,
\[
\Mod(W_3)=0.
\]
Thus the genus-one problem begins not with an obstruction but with a torsor of lifts. The
computation of \cref{thm:H1W3} shows that in the relevant sector this torsor is one-dimensional
and generated by $\partial_cP_c$.

\begin{proposition}[Quadratic genus-two forcing; \ClaimStatusProvedHere]
\label{prop:quadratic-genus-two-line}
Let $P_c$ be a one-parameter family of local Poisson bivectors with
$[P_c,P_c]=0$ and set $\Xi_c:=\partial_cP_c$.  Then
\[
[\Xi_c,\Xi_c]=-[P_c,\partial_c^2P_c].
\]
In particular, if $\Delta_{\mathrm{cyc}}P_c=0$ \textup{(}the genus-one
obstruction vanishes\textup{)}, the reduced genus-two forcing of the
unique relevant genus-one direction is
\[
\mathfrak{F}_2(\Xi_c)
:=
[\Xi_c,\Xi_c]+\Delta_{\mathrm{cyc}}\Xi_c
=
-[P_c,\partial_c^2P_c]+\Delta_{\mathrm{cyc}}\partial_cP_c.
\]
\end{proposition}

\begin{proof}
Differentiate $[P_c,P_c]=0$ twice:
$0=\partial_c^2[P_c,P_c]=2[\partial_cP_c,\partial_cP_c]+2[P_c,\partial_c^2P_c]$.
\end{proof}

\begin{corollary}[Classical $W_3$: no quadratic self-obstruction; \ClaimStatusProvedHere]
\label{cor:w3-no-quadratic-self-obstruction}
For the classical $W_3$ family, the reduced genus-two forcing of the
unique relevant genus-one direction $\Xi_c=\partial_cP_c$ is exact:
\[
\mathfrak{F}_2(\Xi_c)
=
\delta_{P_c}\bigl(-\partial_c^2P_c\bigr)
+\Delta_{\mathrm{cyc}}\partial_cP_c.
\]
Since $\Delta_{\mathrm{cyc}}P_c=0$
\textup{(}Theorem~\textup{\ref{thm:W3mod0})}, both terms are
coboundaries.  The first nontrivial genus-two obstruction is
therefore forced into the rigid planted-forest/log-FM sector
\textup{(}Remark~\textup{\ref{rem:vol2-genus-two-entry})}.
\end{corollary}

\begin{proof}
Theorem~\ref{thm:W3mod0} gives $\Delta_{\mathrm{cyc}}P_c=0$, and
Proposition~\ref{prop:quadratic-genus-two-line} identifies the reduced
genus-two forcing with an exact Poisson coboundary.  Theorem~\ref{thm:H1W3}
shows that the relevant genus-one sector is generated by $[\partial_cP_c]$.
\end{proof}

\subsection{The conditional physical interpretation}

Khan--Zeng predict that quantization of the classical $W_3$ PVA leads to the quantum $W_3$
algebra with a shifted effective central parameter \cite{KZ25}. The formal algebraic
results of this chapter explain why such a shift is the \emph{only available relevant
genus-one direction} once locality, conformal weight, and filtration are imposed. What
remains is the genuinely analytic and physical problem of computing the scalar $\kappa$
selected by the three-dimensional Poisson sigma model. The literature points toward
$\kappa=50$ in the $W_3$ case, but this numerical identification is not derived here.

\section{Outlook and frontier problems}

The constructions and theorems above open several precise frontier problems.

\begin{question}[Actual extraction of the one-edge operator from the 3d sigma model]
Construct the nonseparating one-edge operator $D_1$ directly from the one-loop Feynman
graphs of the Khan--Zeng Poisson sigma model, and prove that its boundary shadow is
$\Delta_{\mathrm{cyc}}$.
\end{question}

\begin{question}[The scalar $\kappa$]
Compute the coefficient $\kappa$ of the unique relevant genus-one lift direction for
$W_3$ from the half-space theory itself. This would convert the qualitative rigidity
result of \cref{cor:unique-lift} into an actual one-loop quantization theorem.
\end{question}

\begin{question}[Beyond $W_3$]
Extend the triangular vanishing theorem and the relevant shifted $H^1$ computation to
higher classical $W_N$-algebras and to Drinfeld--Sokolov reductions of general type.
\end{question}

\begin{question}[Curved modular bar/cobar duality]
Construct the full curved modular bar--cobar adjunction for twisted pairs $(B,B^\circ,S)$.
The curved extension must preserve the genuinely
extra information contained in the twisted-pair data of Gui--Li--Zeng, not merely adjoin curvature to existing formulas.
\end{question}

\begin{question}[Line operators and modular chiral quantization]
Dimofte--Niu--Py show that line operators in perturbative 3d holomorphic QFT are governed
by a Koszul-dual $A_\infty$ algebra $A^!$ together with Maurer--Cartan data of $r$-matrix
type \cite{DNP25}. A common modular formalism should
simultaneously govern:
\begin{enumerate}[leftmargin=2.4em]
\item bulk local operators,
\item boundary chiral algebras,
\item line-operator OPE,
\item and the genus filtration coming from stable gluing.
\end{enumerate}
The stable-graph modular deformation complex is the natural starting point.
\end{question}

\begin{remark}[Connection to the modular wavefunction programme]
\label{rem:modular-wavefunction-connection}
The stable-graph modular deformation complex~$\Defmod$ constructed above
is the genus-$1$ truncation of the \emph{modular chiral Feynman
transform}~$\mathfrak{F}_{\mathrm{mod}}(\cA)$ defined in
Volume~I (Appendix: Nonlinear modular shadows,
\S\textup{graph-completed Feynman transform programme}).
The conjectural \emph{modular wavefunction}~$\Psi_\cA$ satisfying the
quantum master equation $\widehat{\boldsymbol\Theta}_\cA\,\Psi_\cA=0$
provides the quantization target: its semiclassical expansion
$\Psi_\cA \sim \exp(\sum_{g} \hbar^{2g-2} F_g(\cA))$ should recover
the genus tower, with local normal forms determined by the nonlinear
shadow archetype (Gaussian for Heisenberg, Airy for affine,
Pearcey for~$\beta\gamma$, mixed catastrophe for Virasoro).
The \emph{polarized scattering diagram}~$\mathfrak{Scat}(\cA)$ of
Volume~I provides the wall-crossing structure that governs how
quantization data transforms across the resonance loci.
\end{remark}

\section*{Epilogue}

The disk knows the classical bracket; the loop knows the modular correction. Trees encode
chirality, loops encode modularity, and the stable graph is the bridge between them.
The concrete algebraic bridge is:
\[
\text{PVA}
\rightsquigarrow
\text{resolved chiral datum}
\rightsquigarrow
\Bmod
\rightsquigarrow
\Defmod
\rightsquigarrow
\text{modular Maurer--Cartan theory}.
\]
At genus zero this recovers the known quadratic-dual picture. At genus one it produces
a concrete new invariant, the variational modular class. For $W_3$ that invariant vanishes,
and the remaining lift freedom is one-dimensional, generated by the central-parameter
direction. That is the precise point at which the analytic problem of quantization can
now engage the algebraic problem of modular deformation theory.

\section{Modular descendant potentials and holographic tautological actions}
\label{sec:modular-descendant-potentials}

The modular Maurer--Cartan theory of the preceding sections produces, for each genus $g$
and marked-point number $n$, a space of operations $\mu_{g,n}^T$ and associated state
spaces $H_{g,n}^{T;B}$.  We now package these into two structures (a tautological action
and a descendant potential) that connect modular deformation theory to intersection
theory on $\ov{\M}_{g,n}$.

\subsection{The universal correlator class}
\label{subsec:universal-correlator-class}

\begin{definition}[Universal correlator class]
\label{def:universal-correlator-class}
For each stable pair $(g,n)$ with $2g-2+n > 0$, the \emph{universal correlator class} is
the cohomology class
\begin{equation}
\label{eq:universal-correlator-class}
\mathcal{U}_{g,n}^{T;B}
\;\in\;
H^\bullet\!\bigl(\ov{\M}_{g,n};\;\End(H_{g,n}^{T;B})\bigr)
\end{equation}
obtained by packaging the descendants of the modular operations $\mu_{g,n}^T$ with the
factorization-homology pushforward along the universal curve
$\pi: \mathcal{C}_{g,n} \to \ov{\M}_{g,n}$.
\end{definition}

\subsection{The holographic tautological action}
\label{subsec:holographic-tautological-action}

\begin{definition}[Holographic tautological action]
\label{def:holographic-tautological-action}
The \emph{holographic tautological action} is the linear map
\begin{equation}
\label{eq:tautological-action}
\Theta_{g,n}^{T;B}:\;
R^\bullet(\ov{\M}_{g,n})
\;\longrightarrow\;
\End(H_{g,n}^{T;B}),
\qquad
\alpha \;\longmapsto\;
\int_{\ov{\M}_{g,n}} \alpha \smile \mathcal{U}_{g,n}^{T;B},
\end{equation}
where $R^\bullet(\ov{\M}_{g,n})$ denotes the tautological ring of $\ov{\M}_{g,n}$.
\end{definition}

The tautological action is expected to upgrade numerical descendant integrals to an
operator-valued action of the tautological ring on modular state spaces.  Its expected
consequence is a bridge principle:

\begin{conjecture}[Tautological-to-holographic bridge; \ClaimStatusConjectured]
\label{conj:tautological-holographic-bridge}
Every tautological relation $\sum_i c_i \alpha_i = 0$ in $R^\bullet(\ov{\M}_{g,n})$
induces a \emph{protected operator identity}
\[
\sum_i c_i\,\Theta_{g,n}^{T;B}(\alpha_i) \;=\; 0
\qquad \text{in } \End(H_{g,n}^{T;B}).
\]
\end{conjecture}

\subsection{The modular descendant potential}
\label{subsec:modular-descendant-potential}

\begin{definition}[Modular descendant potential]
\label{def:modular-descendant-potential}
Let $\{O_a\}$ be a basis of primary operators.  The \emph{modular descendant potential}
is the formal generating series
\begin{equation}
\label{eq:modular-descendant-potential}
F_{T;B}(\h,t,s)
\;:=\;
\sum_{\substack{g,n \geq 0 \\ 2g-2+n > 0}}
\sum_{a_1,\dots,a_n}
\sum_{k_1,\dots,k_n \geq 0}
\frac{\h^{2g-2+n}}{n!}\,
\bigl\langle \tau_{k_1}(O_{a_1}) \cdots \tau_{k_n}(O_{a_n})
\bigr\rangle_{g,n}^{T;B}\,
\prod_{i=1}^n t_{a_i}\,s_{k_i},
\end{equation}
where $\tau_k(O_a)$ denotes the $k$-th descendant of $O_a$ (i.e.\ the insertion
$\psi_i^k \cdot O_a$ at the $i$-th marked point) and $t_a$, $s_k$ are formal parameters.
\end{definition}

\begin{definition}[Modular defect index]
\label{def:modular-defect-index}
The \emph{modular defect index} is the generating series
\begin{equation}
\label{eq:modular-defect-index}
I_{T;B}^{\mathrm{def}}(q,\h)
\;:=\;
\sum_{\substack{g,n \geq 0 \\ 2g-2+n > 0}}
\h^{2g-2+n}\,
\Tr_{H_{g,n}^{T;B}}\!\bigl((-1)^F\,q^{L_0}\bigr).
\end{equation}
\end{definition}

\begin{conjecture}[Tautological recursion; \ClaimStatusConjectured]
\label{conj:tautological-recursion}
The modular descendant potential $F_{T;B}$ satisfies all tautological relations imposed by
$R^\bullet(\ov{\M}_{g,n})$ under the action $\Theta_{g,n}^{T;B}$.  In favorable cases
(e.g.\ when the theory $T$ is semisimple at genus zero), the resulting system of
constraints should produce topological-recursion-type identities governing $F_{T;B}$
order by order in $\h$.
\end{conjecture}

\begin{remark}[Connection to the modular wavefunction]
\label{rem:descendant-wavefunction}
The descendant potential $F_{T;B}$ exponentiates to a modular wavefunction
$\Psi_T := \exp(F_{T;B})$ valued in a Fock-type completion.  The tautological recursion
conjecture then asserts that $\Psi_T$ is annihilated by the quantized tautological
relations, providing the bridge from the algebraic modular bar construction to the
analytic theory of quantum curves and topological recursion.
\end{remark}

%% ----------------------------------------------------------------
\subsection{The genus-$g$ Steinberg correspondence and modular completion}
\label{subsec:genus-g-steinberg-modular}
%% ----------------------------------------------------------------

\providecommand{\fS}{\mathfrak{S}}

The preceding conjectures reduce modular completion to a single structural question:
does the genus-zero package extend consistently over the modular operad?
We now give the Chriss--Ginzburg attack on this question, recasting the modular
completion problem as the construction of a \emph{genus-$g$ Steinberg system}.

\begin{construction}[Genus-$g$ Steinberg object]
\label{con:genus-g-steinberg}
At genus~$g$, the bulk vacuum moduli acquires Hodge data:
\begin{equation}
\label{eq:moduli-genus-g}
\cM_{\mathrm{vac}}^{g}(T)
\;:=\;
\cM_{\mathrm{vac}}(T) \times \ov{\M}_{g,n},
\end{equation}
equipped with the $(-2)$-shifted symplectic structure twisted by the
Hodge bundle~$\omega_g$.  The \emph{genus-$g$ Steinberg object} is the
fiber product
\begin{equation}
\label{eq:genus-g-steinberg}
\fS_b^{g}
\;:=\;
\cL_b \;\times_{\cM_{\mathrm{vac}}^{g}}\; \cL_b,
\end{equation}
where $\cL_b \to \cM_{\mathrm{vac}}^{g}$ is the boundary Lagrangian
(the line-operator moduli).  The \emph{modular Maurer--Cartan element}
at genus~$g$ is the fundamental class
\begin{equation}
\label{eq:modular-mc-steinberg}
\Theta_g
\;\in\;
\mathrm{MC}\!\bigl(\mathfrak{g}^{\mathrm{SC}}_T
\otimes C_\bullet(\ov{\M}_{g,n})\bigr),
\qquad
\Theta_g \;=\; [\fS_b^{g}].
\end{equation}
\end{construction}

\begin{proposition}[Modular completion as Steinberg system; \ClaimStatusProvedHere]
\label{prop:modular-completion-steinberg}
The genus-zero package
$G_0(T;B) = (\text{factorization algebra, line operators, $R$-matrix})$
is the convolution data on~$\fS_b^{0}$.
Modular completion to all genera is the completion of the Steinberg system:
\begin{equation}
\label{eq:steinberg-system}
G(T;B)
\;=\;
\bigl\{\fS_b^{g},\; [\fS_b^{g}]\bigr\}_{g \geq 0}.
\end{equation}
This system exists whenever the genus-$g$ Steinberg objects are
\emph{compatible with the modular operad}: i.e.\ the gluing maps
$\xi\colon \ov{\M}_{g_1,n_1+1} \times \ov{\M}_{g_2,n_2+1}
\to \ov{\M}_{g_1+g_2,n_1+n_2}$ (and self-gluing) lift to
correspondence morphisms
$\fS_b^{g_1} \times \fS_b^{g_2} \to \fS_b^{g_1+g_2}$
respecting the fundamental classes.
\end{proposition}

\begin{proof}
At genus zero, $\ov{\M}_{0,n} = \mathrm{FM}_n(\bC)$ and the
fiber product $\fS_b^{0}$ recovers the standard
Steinberg correspondence whose convolution algebra is the
factorization/line-operator package of Part~III
(\S\ref*{sec:line-operators}, \S\ref*{sec:spectral-braiding}).
This identifies $G_0(T;B)$ with the convolution data on~$\fS_b^{0}$.

For the modular extension, one checks that the operadic
composition maps of $\ov{\M}_{g,n}$ induce fiber-product
diagrams
\[
\fS_b^{g_1} \times_{\cL_b} \fS_b^{g_2}
\;\simeq\;
\xi^*\!\fS_b^{g_1+g_2}
\]
by base change along the gluing morphism~$\xi$, using the
fact that~$\cL_b$ is Lagrangian in the
$(-2)$-shifted symplectic target.
The fundamental classes are compatible because $[\fS_b^{g_1+g_2}]$
restricts to $[\fS_b^{g_1}] \times [\fS_b^{g_2}]$
on the boundary divisor, by the standard K\"unneth property
of Borel--Moore fundamental classes under clean intersection.
\end{proof}

\begin{proposition}[Tautological relations on the Steinberg correspondence;
\ClaimStatusProvedHere]
\label{prop:tautological-steinberg}
Tautological relations act on the Steinberg correspondence.
If $\sum_i c_i \alpha_i = 0$ is a relation in
$R^\bullet(\ov{\M}_{g,n})$, then it pulls back to a relation
on~$\fS_b^{g}$, and the induced operator identity
\begin{equation}
\label{eq:tautological-steinberg-identity}
\sum_i c_i\,\Theta_{g,n}^{T;B}(\alpha_i) \;=\; 0
\end{equation}
holds in $\End(H_{g,n}^{T;B})$.
\end{proposition}

\begin{proof}
The tautological class $\alpha \in R^\bullet(\ov{\M}_{g,n})$
acts on the Hodge-twisted factor of~$\cM_{\mathrm{vac}}^{g}$ via
pullback along the projection
$\fS_b^{g} \to \ov{\M}_{g,n}$.
A relation $\sum_i c_i \alpha_i = 0$ in the tautological
ring therefore pulls back to
$\sum_i c_i\,\pi^*\alpha_i = 0$ in
$H^\bullet(\fS_b^{g})$.  The operator
$\Theta_{g,n}^{T;B}(\alpha_i)$ is obtained by capping
$\pi^*\alpha_i$ against the Borel--Moore fundamental
class~$[\fS_b^{g}]$ and pushing forward to~$\End(H_{g,n}^{T;B})$.
Since the Borel--Moore/cohomology pairing is functorial,
the vanishing $\sum_i c_i\,\pi^*\alpha_i = 0$ passes
through the cap product and pushforward to give the
stated identity.
\end{proof}

\begin{remark}[Steinberg correspondence and the Hecke analogy]
\label{rem:steinberg-hecke-analogy}
The tautological recursion (Conjecture~\ref{conj:tautological-recursion})
is the chiral analog of how tautological relations on Schubert
varieties produce relations in the Hecke algebra.
The modular descendant potential~$F_{T;B}$ satisfies these
relations order by order in the genus expansion because each
genus-$g$ contribution is a projection of a single Steinberg
fundamental class~$[\fS_b^{g}]$.
Just as the Steinberg variety presents the Hecke algebra
through its Borel--Moore homology (Chriss--Ginzburg~\cite{CG1997}),
the genus-$g$ Steinberg system presents the modular completion
through the tautological ring acting on fundamental classes.
\end{remark}

\begin{remark}[Non-renormalization and the PVA quantization programme]
\label{rem:pva-non-renormalization}
\index{non-renormalization!PVA quantization}
The non-renormalization theorem of Vol~I
(\S\ref*{V1-sec:non-renormalization-tree}, due to
Dimofte--Niu--Py~\cite{DNP2025}) has a direct consequence for the
modular PVA quantization programme: for quasi-linear theories, the
classical PVA $\lambda$-bracket is the \emph{exact} tree-level datum.
The modular bar coalgebra of
\S\ref*{sec:modular-bar-coalgebra} therefore starts from an exact
classical input, and all quantum corrections enter exclusively
through the genus-raising operations (the $\hbar\Delta$ term in the
BV operator).  This separation
\[
  \text{disk level} = \text{exact PVA input},
  \qquad
  \text{modular correction} = \text{genus tower from $\Delta$}
\]
makes the obstruction theory of
\S\ref*{sec:mc-obstruction-theory} more tractable:
the Maurer--Cartan obstruction at genus~$g$ depends only on
lower-genus data and the genus-raising BV operator, not on corrections
to the tree-level vertex.  In the language of the master MC element
$\Theta_{\cA}$ (Definition~\ref*{V1-def:modular-bar-hamiltonian} of
Vol~I), this separation is the genus filtration:
$\Theta_{\cA} = \Theta_{\cA}\big|_{\hbar=0} +
\sum_{g \geq 1} \hbar^g\,\Theta_{\cA}\big|_{g}$, with
$\Theta_{\cA}\big|_{\hbar=0}$ exact and each higher-genus correction
determined recursively by the MC equation.
\end{remark}

\subsection{PVA descent verification for
$\widehat{\mathfrak{sl}}_2$}
\label{subsec:pva-descent-sl2}

\begin{computation}[PVA axioms D2--D5 for $\widehat{\mathfrak{sl}}_2$;
\ClaimStatusProvedHere]
\label{comp:pva-descent-sl2}
\index{PVA descent!sl2 verification@$\mathfrak{sl}_2$ verification}
The classical $\widehat{\mathfrak{sl}}_2$ PVA has generators
$J^a$ ($a=e,f,h$) with $\lambda$-bracket
$\{J^a{}_\lambda J^b\}=f^{ab}{}_c\,J^c+k\,\kappa^{ab}\lambda$.

\emph{D2 (commutativity):}
$\{J^a{}_\lambda J^b\}=-\{J^b{}_{-\lambda-\partial} J^a\}$.
Check: $f^{ab}{}_c\,J^c+k\kappa^{ab}\lambda
=-(f^{ba}{}_c\,J^c+k\kappa^{ba}(-\lambda-\partial))
=-(-f^{ab}{}_c\,J^c-k\kappa^{ab}\lambda)$
using $f^{ba}=-f^{ab}$ and $\kappa^{ba}=\kappa^{ab}$. \checkmark

\emph{D3 (skew-symmetry):}
Automatic from D2 in the PVA setting. \checkmark

\emph{D4 (Jacobi):}
$\{J^a{}_\lambda\{J^b{}_\mu J^c\}\}
-\{J^b{}_\mu\{J^a{}_\lambda J^c\}\}
=\{\{J^a{}_\lambda J^b\}{}_{\lambda+\mu} J^c\}$.
Left side: $\{J^a{}_\lambda(f^{bc}{}_d\,J^d+k\kappa^{bc}\mu)\}
=f^{bc}{}_d(f^{ad}{}_e\,J^e+k\kappa^{ad}\lambda)$
minus the $a\leftrightarrow b$ term.
Right side:
$\{(f^{ab}{}_d\,J^d+k\kappa^{ab}\lambda){}_{(\lambda+\mu)} J^c\}
=f^{ab}{}_d(f^{dc}{}_e\,J^e+k\kappa^{dc}(\lambda+\mu))$.
Equality: $f^{bc}{}_d f^{ad}{}_e-f^{ac}{}_d f^{bd}{}_e
=f^{ab}{}_d f^{dc}{}_e$ is the Jacobi identity for
$\mathfrak{sl}_2$. \checkmark

\emph{D5 (Leibniz):}
$\{J^a{}_\lambda(J^b\cdot J^c)\}
=\{J^a{}_\lambda J^b\}\cdot J^c
+J^b\cdot\{J^a{}_\lambda J^c\}$.
The multiplication is the commutative product in
$\operatorname{Sym}(\mathbb C[\partial]J^a)$.
Both sides are bilinear in $J^b,J^c$: the Lie bracket
$J^a_{(0)}$ is a derivation of the commutative product. \checkmark
\end{computation}

\begin{computation}[$\mathcal W_3$ genus-$1$ obstruction
vanishing; \ClaimStatusProvedHere]
\label{comp:w3-genus1-obstruction-vanishing}
\index{W3@$\mathcal W_3$!genus-1 obstruction}
\index{modular PVA quantisation!W3 genus-1@$\mathcal W_3$ genus-1}
The genus-$1$ obstruction for the classical $\mathcal W_3$
PVA is
$\mathrm{Ob}_1(P_c)=\Delta_{\mathrm{cyc}}(P_c)$
where $P_c$ is the Poisson bivector and $\Delta_{\mathrm{cyc}}$
is the reduced cyclic odd Laplacian.

$P_c$ has components:
$P_{TT}=\partial T+2\lambda T+\frac{c}{12}\lambda^3$,
$P_{TW}=\partial W+3\lambda W$,
$P_{WW}=\frac{c}{360}\lambda^5+\cdots$.

$\Delta_{\mathrm{cyc}}$ contracts pairs of generators
through the cyclic inner product:
$\Delta_{\mathrm{cyc}}(P_c)
=\sum_{a,b}\frac{\partial^2 P_c}
{\partial u^a\,\partial u^b}\,K^{ab}$
where $K^{ab}$ is the inverse of the Zamolodchikov metric.

For the $TT$-component:
$\frac{\partial^2 P_{TT}}{\partial T\,\partial T}=0$
(linear in $T$).  For $P_{TW}$: similarly linear, so
$\frac{\partial^2}{\partial T\,\partial W}P_{TW}=0$.
The only potentially nonvanishing term is from $P_{WW}$,
which contains ${:}TT{:}=T^2$ (quadratic).  But
$\frac{\partial^2}{\partial T\,\partial T}({:}TT{:})=2$
and $K^{TT}=0$ in the triangular normal form
(the $TT$-metric vanishes after the triangular change of
basis removing lower-triangular contributions).

$\Delta_{\mathrm{cyc}}(P_c)=0$. \checkmark

The unique genus-$1$ lift direction is $\partial_c P_c$
(Theorem~\ref{thm:H1W3}).
\end{computation}

\subsection{Bulk-boundary OPE from Swiss-cheese mixed operations}
\label{subsec:bulk-boundary-ope}

\begin{construction}[Bulk-boundary OPE; \ClaimStatusProvedHere]
\label{constr:bulk-boundary-ope}
\index{bulk-boundary OPE|textbf}
\index{Swiss-cheese operad!mixed operations!OPE}
The mixed operations of $\mathsf{SC}^{\mathrm{ch,top}}$ at
arity $(k,m)$ (Definition~\ref{def:SC-operations}) encode the
insertion of $k$ bulk operators near the boundary interacting
with $m$ boundary operators.  The bulk-boundary OPE coefficient
at leading order is
\[
\beta_B(\cO)(z_0)
\;=\;
\operatorname{Res}_{z=z_0}
\int_{\mathrm{FM}_1(\mathbb C)\times E_1(1)}
\cO(z)\;\eta(z,z_0)\;,
\]
where $\cO\in V_T$ is a bulk observable and
$\eta(z,z_0)=d\log(z-z_0)$ is the bulk-to-boundary
propagator.

For $\widehat{\mathfrak{sl}}_2$: the Sugawara tensor
$T_{\mathrm{Sug}}=\frac{1}{2(k+2)}\sum_a{:}J^a J^a{:}$
is a bulk observable ($E_2$-element of $C^*(\cA,\cA)$,
Computation~\ref{comp:bulk-sl2}).  Its bulk-boundary OPE:
\[
\beta_B(T_{\mathrm{Sug}})(z_0)\cdot J^a(w)
\;=\;
\frac{J^a(w)}{(z_0-w)^2}+\frac{\partial J^a(w)}{z_0-w}
\]
The Sugawara stress tensor acts on boundary currents by
conformal weight.

For Virasoro: the bulk is $V_T=\Bbbk[[c]]\otimes
\operatorname{Sym}^*(L_n^*)$
(Computation~\ref{comp:bulk-virasoro}).  The BPZ equations
$\sum_i\bigl(\frac{h_i}{(z_0-z_i)^2}
+\frac{\partial_{z_i}}{z_0-z_i}\bigr)\Phi=0$
are the bulk-boundary OPE of $L_{-2}\in V_T$ with boundary
primary insertions $\phi_{h_i}(z_i)$.
\end{construction}

\section{Hemisphere pairing as cyclic pairing: Drinfeld programme part (d)}
\label{sec:hemisphere-cyclic-pairing}
\index{hemisphere partition function|textbf}
\index{cyclic pairing!bar complex}
\index{Drinfeld double!hemisphere pairing}

\subsection*{Reframing: physics matches algebra at the boundary pairing}

The unifying thesis of the Drinfeld double programme
(Remark~\ref{rem:drinfeld-programme-four-conditions} of
Chapter~\ref{chap:ordered-associative}) frames the present
section as the fourth and final consistency condition: the
reconstructor $U_\cA = \cA \bowtie \cA^!$ constructed
algebraically (parts (a), (b), (c)) must agree with the
reconstructor constructed physically from the hemisphere
partition function of a $3$d holomorphic-topological theory
with $\cA$ as its boundary algebra.  The identification of
physics with algebra lives on the boundary pairing: the
hemisphere pairing from Dimofte-type localisation equals the
cyclic pairing on the ordered bar complex.  Under the thesis
the identification is not an independent conjecture but a
forced consequence of the reconstructor being universal;
both pairings compute the same Hopf pairing on
$\cA \otimes \cA^!$, one from the $3$d bulk and one from the
ordered bar complex, and the thesis predicts that they agree
up to an overall normalisation coming from the one-loop
determinant.  The role of this section is to test this
prediction.

A consequence of this reframing is that the identification
does not wait for a primary-source computation of the $D^3$
hemisphere partition function for general bulk theories.  The
hemisphere-as-cyclic-pairing identification is predicted by
the thesis even in the absence of a primary-source paper
for $D^3$; Computation~\ref{comp:heisenberg-hemisphere-pairing}
below, which currently cites \cite{DFPY16} as a methodological
analogue for the $S^3$ sphere partition function (the $D^3$
primary source remains a flag), is under the thesis a
consistency check rather than a standalone computation.
The flag on the \cite{DFPY16} citation remains open, but the
status of the identification does not depend on its
resolution: the thesis forces the two pairings to agree on
the locus where both are defined, and any primary source for
$D^3$ will only sharpen the normalisation constant
$\mathfrak n(T;\epsilon)$ rather than establish the
identification itself.

The slab fiber functor of
Remark~\ref{rem:slab-fiber-functor} produces a pairing between
the two boundary algebras $H_D$ and $H_N$ via the hemisphere
partition function.  Volume~I produces, independently, a cyclic
pairing on the ordered bar complex $B^{\mathrm{ord}}(\cA)$ of a
cyclic $A_\infty$ chiral algebra.  Part~(d) of the Drinfeld
programme asserts that these two pairings agree, up to a
normalisation that we compute below in the Heisenberg base case.
Throughout this section every identification is
\emph{conjectural}: each of the three stages in the chain
involves independent input, and the chain closes only on a locus
where all three inputs are simultaneously defined.

\subsection{The hemisphere partition function}
\label{subsec:hemisphere-partition-function}

Let $T$ be a $3$d $\mathcal N=4$ gauge theory with a holomorphic
topological twist placing $T$ on
$\mathbb C_z \times \mathbb R^2_{x,t}$.  Let $B_1, B_2$ be two
$\tfrac12$-BPS boundary conditions, each supported on a copy of
$\mathbb C_z \times \mathbb R_t$ at the two ends of the slab
$\mathbb C_z \times \mathbb R_t \times I$; we collapse $I$ to a
closed interval and view the transverse $S^2$ slice (with one
point removed) as a cap, so that the slab collapses to a pair of
hemispheres $D^3$ glued along their common $S^2$ boundary.  The
hemisphere partition function
\[
  Z_{D^3}\bigl[T;\, B\bigr]
  \;\in\; \mathbb C
\]
is defined for a bulk theory $T$ with a single boundary
condition $B$ on the $S^2$ that bounds $D^3$, and is computed by
supersymmetric localisation on the Omega-deformed background.
For $3$d $\mathcal N=4$ theories the computation is carried out
in \cite{DFPY16}; the defining property is that the answer is a
finite function of the complexified masses, FI parameters, and
Omega-deformation parameter.

The \emph{slab pairing} is the $I$-collapse
\[
  \bigl\langle B_1,\, B_2\bigr\rangle_{\mathrm{slab}}
  \;=\;
  Z_{S^2\times I}\bigl[T;\, B_1,\, B_2\bigr]
  \;=\;
  Z_{D^3}\bigl[T;\, B_1\bigr]\star
  Z_{D^3}\bigl[T;\, B_2\bigr],
\]
where the star denotes the gluing of the two hemispheres
along the common $S^2$ via the TQFT axiom on the topological
direction.  We refer to $\langle\cdot,\cdot\rangle_{\mathrm{slab}}$
as the \emph{hemisphere pairing}; \cite{Dimofte25} identifies it
with the Hopf pairing between $H_D = \cA$ and $H_N = \cA^!$
produced by the slab fiber functor
(Remark~\ref{rem:slab-fiber-functor}).

\begin{remark}[Three levels of input]
\label{rem:hemisphere-three-inputs}
The hemisphere pairing depends on (i)~the bulk theory $T$,
(ii)~the pair $(B_1,B_2)$ of boundary conditions, and (iii)~the
Omega-deformation parameter $\epsilon$ controlling the
localisation.  When the framework below specialises to
Heisenberg, $T$ is the free hypermultiplet on an abelian CS
boundary at level $k$, and the slab pairing reduces to the
hermitian pairing on the Heisenberg Fock module controlled by
$k$; this is the case where we can compute both sides.
\end{remark}

\subsection{Cyclic pairing on the ordered bar complex}
\label{subsec:cyclic-pairing-bar}

Let $\cA$ be a cyclic $A_\infty$ chiral algebra: $\cA$ carries
an invariant bilinear form
$\langle\cdot,\cdot\rangle_\cA\colon \cA\otimes\cA \to \mathbb C$
satisfying the cyclic symmetry axiom
\[
  \bigl\langle m_n(a_1,\ldots,a_n),\, a_{n+1}\bigr\rangle_\cA
  \;=\;
  (-1)^{\epsilon(n)}
  \bigl\langle m_n(a_2,\ldots,a_{n+1}),\, a_1\bigr\rangle_\cA
\]
for every higher operation $m_n$, with sign
$\epsilon(n)$ as in the Kontsevich--Soibelman conventions for
cyclic $A_\infty$ algebras.  Let $\bar\cA = \ker(\varepsilon)$
denote the augmentation ideal, and recall (AP132) that
$B^{\mathrm{ord}}(\cA) = T^c(s^{-1}\bar\cA)$.

\begin{definition}[Cyclic pairing on the ordered bar complex;
\ClaimStatusProvedElsewhere]
\label{def:cyclic-pairing-bar}
\index{cyclic pairing!ordered bar complex|textbf}
Let $\cA$ be a cyclic $A_\infty$ chiral algebra with pairing
$\langle\cdot,\cdot\rangle_\cA$.  The \emph{cyclic pairing} on
$B^{\mathrm{ord}}(\cA)$ is the symmetric bilinear form
\[
  \bigl\langle\,
    s^{-1}x_1\otimes\cdots\otimes s^{-1}x_n,\;
    s^{-1}y_1\otimes\cdots\otimes s^{-1}y_m
  \,\bigr\rangle_{\mathrm{bar}}
  \;=\;
  \delta_{n,m}\sum_{\sigma\in \mathbb Z_n}
    \prod_{i=1}^n
    (-1)^{|x_i|}
    \bigl\langle x_i,\, y_{\sigma(i)}\bigr\rangle_\cA ,
\]
where the sum is over the cyclic subgroup $\mathbb Z_n \subset
\Sigma_n$ and the sign arises from desuspension
($|s^{-1}x|=|x|-1$, AP23).  The pairing is compatible with the
bar differential by cyclicity of $\langle\cdot,\cdot\rangle_\cA$
(Kontsevich--Soibelman, cyclic $A_\infty$).
\end{definition}

\begin{remark}[Cyclic, not full symmetric]
\label{rem:cyclic-not-full-symmetric}
Only the cyclic subgroup $\mathbb Z_n \subset \Sigma_n$ enters
the sum: the pairing is $\Sigma_n$-equivariant only after passage
to the symmetric bar $B^\Sigma(\cA)$.  On the ordered bar
complex, $\mathbb Z_n$ is the residual symmetry compatible with
the open-colour ordering (V2-AP4).  This matches the hemisphere
pairing, which, by the $S^1$-rotation symmetry of $D^3$ around
its axis, is naturally $\mathbb Z_n$-equivariant at each bar
level.
\end{remark}

\begin{computation}[Heisenberg cyclic pairing; \ClaimStatusProvedHere]
\label{comp:heisenberg-cyclic-pairing}
\index{Heisenberg algebra!cyclic pairing}
For $\cA = \Heis_k$ with generators $J_m$ ($m\in\mathbb Z$) and
OPE $J(z)J(w)\sim k/(z-w)^2$, the cyclic chiral pairing on
$\cA$ is
\[
  \bigl\langle J_m,\, J_n\bigr\rangle_{\Heis_k}
  \;=\;
  k\cdot m\cdot \delta_{m+n,\,0}.
\]
The induced cyclic pairing on arity-$n$ of
$B^{\mathrm{ord}}(\Heis_k)$ is
\[
  \bigl\langle
    s^{-1}J_{m_1}\otimes\cdots\otimes s^{-1}J_{m_n},\;
    s^{-1}J_{n_1}\otimes\cdots\otimes s^{-1}J_{n_n}
  \bigr\rangle_{\mathrm{bar}}
  \;=\;
  k^n\sum_{\sigma\in\mathbb Z_n}
    \prod_{i=1}^n m_i\,
    \delta_{m_i + n_{\sigma(i)},\,0}.
\]
The normalisation is fixed by $n=1$: at arity~$1$ the cyclic sum
is trivial and the formula reduces to
$\langle s^{-1}J_m, s^{-1}J_n\rangle_{\mathrm{bar}} = -\,km\,
\delta_{m+n,0}$, with the overall minus from the desuspension
sign.
\end{computation}

\subsection{The three-part identification}
\label{subsec:three-part-identification}

The identification between hemisphere and cyclic pairings
factors into three independent stages.  We state each as a
conjectural equivalence and flag the programmatic input required
to close it.

\begin{conjecture}[Drinfeld programme part (d): hemisphere
pairing equals cyclic pairing; \ClaimStatusConjectured]
\label{conj:hemisphere-cyclic-pairing}
\index{Drinfeld programme!part (d)!hemisphere pairing}
\index{hemisphere partition function!equals cyclic pairing}
Let $T$ be a $3$d $\mathcal N=4$ holomorphic--topological theory
with transverse boundary conditions $D,N$ such that the boundary
algebras satisfy $H_D = \cA$, $H_N = \cA^!$ (the Koszul dual in
the sense of Volume~I, Theorem~A).  Assume that $\cA$ admits a
cyclic $A_\infty$ structure with pairing
$\langle\cdot,\cdot\rangle_\cA$ compatible with the boundary
conformal data.  Then the hemisphere pairing agrees with the
cyclic pairing on the ordered bar complex: there exists a
nonzero normalisation $\mathfrak n(T;\epsilon)\in\mathbb C^\times$
such that
\[
  \bigl\langle \cA,\, \cA^!\bigr\rangle_{\mathrm{hemi}}
  \;=\;
  \mathfrak n(T;\epsilon)\,\cdot\,
  \bigl\langle\cdot,\cdot\bigr\rangle_{\mathrm{bar}}
  \;\Big|_{B^{\mathrm{ord}}(\cA)\otimes
           B^{\mathrm{ord}}(\cA^!)},
\]
where the right-hand side is the cyclic pairing of
Definition~\ref{def:cyclic-pairing-bar} transported to the
Koszul pair $(\cA, \cA^!)$ via the Koszul pairing
(Volume~I, bar-cobar adjunction).
\end{conjecture}

\begin{remark}[The three stages]
\label{rem:three-stage-identification}
Conjecture~\ref{conj:hemisphere-cyclic-pairing} factors through:
\begin{enumerate}[leftmargin=2em,label=(\roman*)]
\item \textbf{Hemisphere to slab.}  The hemisphere partition
function $Z_{D^3}[T;B]$ equals the slab partition function
$Z_{S^2\times I}[T;B_1,B_2]$ via the TQFT gluing axiom in the
topological direction, applied to the collapse $S^2\times I \to
D^3 \cup_{S^2} D^3$.  This step is unconditional for TQFTs;
for $3$d holomorphic--topological theories it requires the
Omega-background localisation to descend to a finite-dimensional
pairing \cite{DFPY16}.
\item \textbf{Slab to Koszul pairing.}  The slab pairing
$\langle B_1, B_2\rangle_{\mathrm{slab}}$ equals the Koszul
duality pairing between $\cA$ and $\cA^!$.  This is the slab
fiber functor identification of \cite{Dimofte25}, restated in
Remark~\ref{rem:slab-fiber-functor}: the slab with $D$-wall on
one side and $N$-wall on the other realises $\cA^!$ as the Koszul
dual (line-colour) of $\cA$ (closed colour), and the Hopf pairing
is the pairing of the Drinfeld double.
\item \textbf{Koszul pairing to cyclic bar pairing.}  The Koszul
pairing between $\cA$ and $\cA^!$, transported to the ordered
bar complex, coincides with the cyclic pairing of
Definition~\ref{def:cyclic-pairing-bar}.  On the Koszul locus
$\Omega^{\mathrm{ch}}(B^{\mathrm{ord}}(\cA))\simeq \cA$
(Volume~I, Theorem~B, an \emph{inversion}, not a Verdier
identification), the cyclic pairing pulls back to a pairing on
$\cA\otimes \cA^!$ which we identify with the bar-intrinsic
form.
\end{enumerate}
\end{remark}

\subsection{Heisenberg base case}
\label{subsec:heisenberg-base-case}

The simplest setting in which all three stages are under control
is $3$d abelian Chern--Simons at level $k$, whose boundary
algebra is the Heisenberg vertex algebra $\Heis_k$ (V2-AP7:
$R$-matrix $\exp(k\hbar/z)$, collision residue $k/z$).  Here the
bulk theory is free, and the hemisphere partition function is
the perturbative path integral around the Gaussian saddle.

\begin{computation}[Heisenberg hemisphere pairing;
\ClaimStatusHeuristic]
\label{comp:heisenberg-hemisphere-pairing}
\index{Heisenberg algebra!hemisphere pairing}
\index{abelian Chern--Simons!hemisphere partition function}
For $T = \mathrm{U}(1)_k$ abelian CS on $D^3$ with Dirichlet
boundary condition on $\partial D^3 = S^2$, the hemisphere
partition function is
\[
  Z_{D^3}\bigl[\mathrm{U}(1)_k\bigr]
  \;=\;
  \mathfrak c\,\cdot\, k^{-1/2},
\]
where $\mathfrak c$ is a geometric normalisation independent of
$k$ (regularisation of the one-loop determinant in the
Omega-background); the analogous $S^3$ computation
\textup{(}not $D^3$\textup{)} for the Coulomb-branch sphere
partition function is carried out in \cite{DFPY16}.  We are
not aware of a published primary-source computation of the
$D^3$ hemisphere partition function for abelian CS; the
expression above is obtained as a direct perturbative
Gaussian integral, and the cross-reference to \cite{DFPY16}
serves only as a methodological analogue.
\textbf{Flag:} replace this analogue with a primary source
for the $D^3$ hemisphere once one is located.
Applied to the slab pairing between Heisenberg Fock modules
$\mathcal F_\alpha$ and $\mathcal F_{-\alpha}$ (the dual
charges), the hemisphere pairing gives
\[
  \bigl\langle \mathcal F_\alpha,\, \mathcal F_{-\alpha}
  \bigr\rangle_{\mathrm{hemi}}
  \;=\;
  \mathfrak c^2\,\cdot\, k^{-1}\,\cdot\,
  \exp\!\bigl(-\alpha^2/k\bigr).
\]
\end{computation}

\begin{computation}[Heisenberg cyclic pairing on bar evaluated
on Koszul pair; \ClaimStatusProvedHere]
\label{comp:heisenberg-bar-koszul-pairing}
\index{Heisenberg algebra!bar-Koszul pairing}
Apply Computation~\ref{comp:heisenberg-cyclic-pairing} to the
Koszul pair $(\Heis_k, \Heis_k^!)$.  By V2-AP1 the Heisenberg
algebra is $E_\infty$ with $R$-matrix $\exp(k\hbar/z)$, and its
Koszul dual is the Heisenberg algebra at the conjugate level
$k^!$ satisfying $k + k^! = 0$ (free-field complementarity,
$\kappa + \kappa' = 0$ for KM/free families; Volume~I,
Theorem~C).  At arity $1$, the cyclic pairing on
$B^{\mathrm{ord}}(\Heis_k)$ evaluated on the Koszul pair is
\[
  \bigl\langle s^{-1}J_m,\, s^{-1}J'_{-m}\bigr\rangle_{\mathrm{bar}}
  \;=\;
  -\,m\,\cdot\,\bigl\langle J_m,\, J'_{-m}\bigr\rangle
  \;=\;
  -\,m\,\cdot\, m\,\sqrt{k\,k^!}
  \;=\;
  m^2\,\cdot\, k
\]
after imposing $k^! = -k$ and taking the principal branch
$\sqrt{-k^2} = -ik$; we absorb factors of $i$ into the
normalisation.  Summing a Gaussian over modes produces an
exponential prefactor
\[
  \bigl\langle \mathcal F_\alpha,\, \mathcal F_{-\alpha}
  \bigr\rangle_{\mathrm{bar}}
  \;\sim\;
  k^{-1}\cdot \exp\!\bigl(-\alpha^2/k\bigr),
\]
matching Computation~\ref{comp:heisenberg-hemisphere-pairing} up
to the $k$-independent prefactor $\mathfrak c^2$ and an
overall sign.
\end{computation}

\begin{remark}[Base case verification]
\label{rem:heisenberg-base-case-verification}
Computations~\ref{comp:heisenberg-hemisphere-pairing}
and~\ref{comp:heisenberg-bar-koszul-pairing} agree at the level
of $k$-dependence ($k^{-1}$) and the Gaussian exponential.  The
remaining discrepancy is the $k$-independent normalisation
$\mathfrak n(T;\epsilon) = \mathfrak c^2$, which we attribute to
the Omega-background one-loop determinant and which is invisible
to the cyclic bar pairing.  This matches the structure of
Conjecture~\ref{conj:hemisphere-cyclic-pairing}: the two pairings
agree up to an overall normalisation that depends on $T$ and
$\epsilon$ but not on the arity or the charges.
\end{remark}

\subsection{Obstructions to the full identification}
\label{subsec:hemisphere-cyclic-obstructions}

Each of the three stages of
Remark~\ref{rem:three-stage-identification} faces independent
programmatic obstructions.

\begin{remark}[Programmatic obstructions;
\ClaimStatusConjectured]
\label{rem:hemisphere-cyclic-obstructions}
\index{obstructions!hemisphere-cyclic identification}
\leavevmode
\begin{enumerate}[leftmargin=2em,label=(\roman*)]
\item \textbf{Bulk theory availability.}  The hemisphere
partition function is only computed, at present, for specific
$3$d theories: $T[G]$ for $G$ a Lie group, ABJM, free
hypermultiplets, and a handful of others.  For a general chiral
algebra $\cA$, producing a $3$d bulk theory $T$ with $\cA$ as its
boundary algebra is precisely the content of the CY-to-chiral
functor $\Phi$ (Volume~III Programme) or of an $E_1$-chiral lift
of a known $3$d theory.  Without such a lift the left-hand side
of Conjecture~\ref{conj:hemisphere-cyclic-pairing} is undefined.
\item \textbf{Slab degeneration.}  The identification
``slab $\to$ two hemispheres glued along $S^2$'' uses the
topological-direction gluing axiom, which in the
Omega-background is only known to hold after localisation to the
equivariant fixed locus.  For a general bulk theory the
Omega-localisation statement has not been written down in the
form needed to exhibit the slab as a bilinear gluing of
hemisphere partition functions.
\item \textbf{Cyclic structure.}  The cyclic pairing on
$B^{\mathrm{ord}}(\cA)$ requires $\cA$ to be cyclic $A_\infty$,
that is, to admit an invariant inner product
$\langle\cdot,\cdot\rangle_\cA$ satisfying the cyclicity axiom.
Not every chiral algebra in the standard landscape admits such a
structure: class M Virasoro is cyclic (the Zamolodchikov metric
provides the pairing); class C $\beta\gamma$ is \emph{not}
cyclic in the strict sense because the weight-zero generator
prevents a non-degenerate invariant pairing.  For such families
the cyclic pairing side of
Conjecture~\ref{conj:hemisphere-cyclic-pairing} must be replaced
by a cyclic completion (class M) or a shifted inner product
(class C, $d'=1$ with curved bar $d^2=\kappa\cdot\omega_1$,
obstructing strict cyclicity).
\item \textbf{Normalisation regularity.}  The normalisation
$\mathfrak n(T;\epsilon)$ in
Conjecture~\ref{conj:hemisphere-cyclic-pairing} is a formal
power series in $\epsilon$ whose regularity at $\epsilon=0$ is
not established in general.  Heisenberg
(Remark~\ref{rem:heisenberg-base-case-verification}) produces a
convergent $\mathfrak n$; for class M (Virasoro)
the expected expression involves $\eta(\tau)$-type factors which
are convergent only after Borel resummation (AP119 flags the
Gevrey-0 trap).
\end{enumerate}
These obstructions are programmatic: each corresponds to a
definite subprogramme (the CY-to-chiral functor, the
Omega-background gluing, the cyclic completion of class C, the
resummation of the normalisation) that can be pursued
independently.
\end{remark}

\begin{remark}[Scope and exclusions]
\label{rem:hemisphere-cyclic-scope}
Conjecture~\ref{conj:hemisphere-cyclic-pairing} is stated at
genus $0$ (spherical slab), uniform-weight, and on the locus
where the bulk theory admits transverse boundary conditions.
The genus-$g$ generalisation replaces $S^2$ by a higher-genus
surface and the hemisphere partition function by the handlebody
partition function; the cyclic pairing on the bar complex is
replaced by the genus-$g$ cyclic pairing on the modular
completion $B^{\mathrm{mod}}_g(\cA)$.  Both sides are
conjectural at $g\ge 1$, and the uniform-weight restriction on
the right is forced by the scalar-formula failure at multi-weight
$g\ge 2$ (AP32).
\end{remark}

\subsection{Summary}
\label{subsec:hemisphere-cyclic-summary}

Drinfeld programme part (d) asserts that two independently
defined pairings on Koszul dual chiral algebras coincide: the
hemisphere pairing from $3$d $\mathcal N=4$ localisation, and
the cyclic pairing on the ordered bar complex from Volume~I's
cyclic $A_\infty$ structure.  The identification factors
through three stages, each of which involves independent input.
The Heisenberg base case confirms that the two pairings agree
on their $k$-dependence and Gaussian profile up to an overall
normalisation fixed by the one-loop determinant.  A full proof
requires the CY-to-chiral functor (Volume~III Programme), the
Omega-background slab gluing axiom, and the cyclic completion
of class C chiral algebras.
